% =========================================================
% Extremal Elements
% =========================================================

\section{Extremal Elements}
\label{sec:extremal-elements}

\begin{remark*}[Toolkit: Extremal Elements]
A poset may contain elements that sit at the boundary of the order.  There
are two distinct senses of ``extreme'': an element can be extreme
\emph{locally} (nothing in the poset exceeds it, though it may be
incomparable with some elements) or extreme \emph{globally} (it dominates
or is dominated by every element).  The definitions below make these
distinctions precise, proceeding from global extremes of the whole poset,
through local extremes of a subset, to the proposition that finite posets
always contain the local kind.
\end{remark*}

% ---------------------------------------------------------
% def:bottom-element
% ---------------------------------------------------------

\begin{definitionbox}{Definition (Bottom Element)}
\begin{definition}[Bottom Element]
\label{def:bottom-element}
Let $(P, \leq)$ be a partially ordered set.  An element $\bot \in P$ is
a \emph{bottom element} (or \emph{least element}) of $P$ if
\[
\forall x \in P, \quad \bot \leq x.
\]
If a bottom element exists it is unique; it is denoted $\bot_P$ or
simply $\bot$.
\end{definition}
\end{definitionbox}

\begin{remark*}[Standard quantified statement]
\[
\operatorname{Bottom}(\bot, P, \leq)
\;\coloneqq\;
\bot \in P \;\land\; \forall x \in P,\; \bot \leq x.
\]
\end{remark*}

\begin{remark*}[Interpretation]
A bottom element is below every element of the poset without exception.
In $(\mathcal{P}(X), \subseteq)$ the bottom element is $\varnothing$.
The natural numbers have a bottom element; the integers do not.
An antichain with more than one element has no bottom element, since no
member is below any other.
\end{remark*}

\begin{dependencies}
\hyperref[def:lattice-order-partial-order]{Partial order}.
\end{dependencies}

% ---------------------------------------------------------
% def:top-element
% ---------------------------------------------------------

\begin{definitionbox}{Definition (Top Element)}
\begin{definition}[Top Element]
\label{def:top-element}
Let $(P, \leq)$ be a partially ordered set.  An element $\top \in P$ is
a \emph{top element} (or \emph{greatest element}) of $P$ if
\[
\forall x \in P, \quad x \leq \top.
\]
If a top element exists it is unique; it is denoted $\top_P$ or simply
$\top$.
\end{definition}
\end{definitionbox}

\begin{remark*}[Standard quantified statement]
\[
\operatorname{Top}(\top, P, \leq)
\;\coloneqq\;
\top \in P \;\land\; \forall x \in P,\; x \leq \top.
\]
\end{remark*}

\begin{remark*}[Interpretation]
Dual to the bottom element via the
\hyperref[prop:duality-principle]{Duality Principle}.
In $(\mathcal{P}(X), \subseteq)$ the top element is $X$ itself.
\end{remark*}

\begin{dependencies}
\hyperref[def:bottom-element]{Bottom element},
\hyperref[prop:duality-principle]{duality principle}.
\end{dependencies}

% ---------------------------------------------------------
% def:maximal-element
% ---------------------------------------------------------

\begin{definitionbox}{Definition}
\begin{definition}[Maximal Element]
\label{def:maximal-element}
Let $(P, \leq)$ be a partially ordered set and let $Q \subseteq P$.  An
element $a \in Q$ is \emph{maximal in $Q$} if no element of $Q$ strictly
exceeds it:
\[
\forall x \in Q, \quad a \leq x \;\Longrightarrow\; a = x.
\]
The set of maximal elements of $Q$ is denoted $\operatorname{Max}(Q)$.
\end{definition}
\end{definitionbox}

\begin{remark*}[Standard quantified statement]
\[
\operatorname{Maximal}(a, Q, \leq)
\;\coloneqq\;
a \in Q \;\land\; \forall x \in Q,\; a \leq x \Rightarrow a = x.
\]
\end{remark*}

\begin{remark*}[Negated quantified statement]
\[
\neg\operatorname{Maximal}(a, Q, \leq)
\;\coloneqq\;
\exists x \in Q,\; a < x,
\]
that is, some element of $Q$ strictly exceeds $a$.
\end{remark*}

\begin{remark*}[Interpretation]
A maximal element asserts only that nothing in $Q$ strictly dominates it.
It need not be comparable with every element of $Q$, so a poset can have
multiple maximal elements simultaneously.  In domain theory, maximal
elements correspond to fully-specified or total objects.
\end{remark*}

\begin{dependencies}
\hyperref[def:lattice-order-partial-order]{Partial order}.
\end{dependencies}

% ---------------------------------------------------------
% def:minimal-element
% ---------------------------------------------------------

\begin{definitionbox}{Definition}
\begin{definition}[Minimal Element]
\label{def:minimal-element}
Let $(P, \leq)$ be a partially ordered set and let $Q \subseteq P$.  An
element $a \in Q$ is \emph{minimal in $Q$} if no element of $Q$ is
strictly below it:
\[
\forall x \in Q, \quad x \leq a \;\Longrightarrow\; x = a.
\]
\end{definition}
\end{definitionbox}

\begin{remark*}[Standard quantified statement]
\[
\operatorname{Minimal}(a, Q, \leq)
\;\coloneqq\;
a \in Q \;\land\; \forall x \in Q,\; x \leq a \Rightarrow x = a.
\]
\end{remark*}

\begin{remark*}[Interpretation]
Dual to maximal element via the
\hyperref[prop:duality-principle]{Duality Principle}.
\end{remark*}

\begin{dependencies}
\hyperref[def:maximal-element]{Maximal element},
\hyperref[prop:duality-principle]{duality principle}.
\end{dependencies}

% ---------------------------------------------------------
% def:greatest-element
% ---------------------------------------------------------

\begin{definitionbox}{Definition}
\begin{definition}[Greatest Element]
\label{def:greatest-element}
Let $(P, \leq)$ be a partially ordered set and let $Q \subseteq P$.  An
element $m \in Q$ is the \emph{greatest element} (or \emph{maximum}) of
$Q$ if it dominates every element of $Q$:
\[
\forall x \in Q, \quad x \leq m.
\]
If it exists it is unique; it is denoted $\max Q$.
\end{definition}
\end{definitionbox}

\begin{remark*}[Standard quantified statement]
\[
\operatorname{Greatest}(m, Q, \leq)
\;\coloneqq\;
m \in Q \;\land\; \forall x \in Q,\; x \leq m.
\]
\end{remark*}

\begin{remark*}[Interpretation]
A greatest element is a maximal element that is additionally comparable
with every element of $Q$.  When $\max Q$ exists,
$\operatorname{Max}(Q) = \{\max Q\}$, but the converse fails: a unique
maximal element need not dominate every element.
\end{remark*}

\begin{dependencies}
\hyperref[def:maximal-element]{Maximal element}.
\end{dependencies}

% ---------------------------------------------------------
% def:least-element
% ---------------------------------------------------------

\begin{definitionbox}{Definition}
\begin{definition}[Least Element]
\label{def:least-element}
Let $(P, \leq)$ be a partially ordered set and let $Q \subseteq P$.  An
element $m \in Q$ is the \emph{least element} (or \emph{minimum}) of $Q$
if it lies below every element of $Q$:
\[
\forall x \in Q, \quad m \leq x.
\]
If it exists it is unique; it is denoted $\min Q$.
\end{definition}
\end{definitionbox}

\begin{remark*}[Standard quantified statement]
\[
\operatorname{Least}(m, Q, \leq)
\;\coloneqq\;
m \in Q \;\land\; \forall x \in Q,\; m \leq x.
\]
\end{remark*}

\begin{remark*}[Interpretation]
Dual to greatest element.
\end{remark*}

\begin{dependencies}
\hyperref[def:greatest-element]{Greatest element},
\hyperref[prop:duality-principle]{duality principle}.
\end{dependencies}

% ---------------------------------------------------------
% prop:finite-poset-maximal
% ---------------------------------------------------------

\begin{propositionbox}{Proposition}
\begin{proposition}[Every Nonempty Finite Poset Has a Maximal Element]
\label{prop:finite-poset-maximal}
Let $(P, \leq)$ be a finite partially ordered set.  Then every nonempty
subset $Q \subseteq P$ has at least one maximal element.  Furthermore,
for every $x \in P$ there exists $y \in \operatorname{Max}(P)$ with
$x \leq y$.
\hyperref[prf:finite-poset-maximal]{Go to proof.}
\end{proposition}
\end{propositionbox}

\begin{remark*}[Standard quantified statement]
\[
P \text{ finite},\; Q \subseteq P,\; Q \neq \varnothing
\;\Longrightarrow\;
\operatorname{Max}(Q) \neq \varnothing.
\]
\[
\forall x \in P,\; \exists y \in \operatorname{Max}(P),\; x \leq y.
\]
\end{remark*}

\begin{remark*}[Interpretation]
Finiteness is essential.  The collection of all finite subsets of
$\mathbb{N}$, ordered by inclusion, is infinite and has no maximal
element.  In a finite poset one can always ascend from any element to
a ceiling.  This property is the finite shadow of Zorn's Lemma and
underlies ascending chain arguments throughout algebra and order theory.
\end{remark*}

\begin{dependencies}
\hyperref[def:maximal-element]{Maximal element},
\hyperref[def:lattice-order-partial-order]{partial order}.
\end{dependencies}
